\section{Finite dilations and arbitrarily long ordinary moment blocks}
\label{stdil:section}

The scalar transform of the balanced Stieltjes report
\cite{stmell:balanced}, recalled in Lemma~\ref{stbil:scalar}, also permits
deliberately chosen dilations. The shift filters in that report impose
linear conditions on Bernoulli polynomials. Here a single shift is retained
and a finite dilation polynomial cancels selected powers of the tail.
This gives arbitrarily long blocks of ordinary moments, with elementary
Bernoulli values when the zero is supplied by the dilation. It also
completes the principal-character case left outside
Section~\ref{stchar:section}. Mellin scaling and finite polynomial
annihilation supply the elementary mechanism; the zeta evaluations
used below are classical \cite{stchar:classical}.

\subsection{The finite-dilation transform and its scalar subtraction}

For positive scales $h_1,\ldots,h_J$, complex coefficients
$c_1,\ldots,c_J$, and a fixed $a>0$, put
\begin{align}
 D(s)&=\sum_{j=1}^Jc_jh_j^{-s},\qquad
 Z(s)=D(s)\zeta(s,a),\notag\\
 H_q(x)&=\sum_{j=1}^J\frac{c_j}{h_j}C_q(x/h_j,a),
 \qquad q\in\mathbb Z_{\ge0}.\label{stdil:kernel}
\end{align}
Every kernel is an explicit finite combination of generalized Stieltjes
functions by \eqref{stbil:kernel}. The uncolored scalar subtraction is
part of its definition. More precisely, if
\[
 F(s,x)=\sum_jc_jh_j^{-s}\zeta(s,x/h_j+a),
\]
then
\begin{equation}\label{stdil:subtraction}
 H_q(x)=\left.\partial_z^q
       \left[x^zF(1+z,x)-\frac{D(1)}z\right]\right|_{z=0}.
\end{equation}
In general $D(1)\ne0$. The additional logarithmic terms inside the
individual $C_q(x/h_j,a)$ must therefore be retained.

\begin{theorem}[Dilation transform and exact convergence criterion]
\label{stdil:criterion}
For $0<\Re u<q+1$,
\begin{equation}\label{stdil:mellin}
 \int_0^\infty x^{u-1}H_q(x)\,dx
    =\pi\csc(\pi u)Q_q(u)Z(1-u).
\end{equation}
For an integer $m\ge q+1$, the ordinary moment of weight $x^{m-1}$
converges if and only if
\begin{equation}\label{stdil:zeros}
 D(-k)\zeta(-k,a)=0\qquad(q\le k\le m-1).
\end{equation}
When it converges it is absolutely convergent and equals
\begin{equation}\label{stdil:moment}
 \int_0^\infty x^{m-1}H_q(x)\,dx
   =(-1)^{m+1}Q_q(m)Z'(1-m).
\end{equation}
These statements also allow zero coefficients and vanishing $D(1)$.
\end{theorem}
\begin{proof}
Substituting $x=h_jy$ in the absolutely convergent scalar transform
multiplies it by $c_jh_j^{u-1}$; summing proves
\eqref{stdil:mellin}. The scalar Euler--Maclaurin expansion gives
\begin{equation}\label{stdil:tail}
 H_q(x)=\sum_{k=q}^{N}\frac{(-1)^kD(-k)\zeta(-k,a)}{k!}
  \left.\partial_z^q(1+z)_k\right|_{z=0}x^{-k-1}
  +O(x^{-N-2}).
\end{equation}
Here the classical identity
$\zeta(-k,a)=-B_{k+1}(a)/(k+1)$ converts the uncolored Bernoulli
coefficients. For $k\ge q$ the Pochhammer derivative is nonzero.
The first surviving term with $k\le m-1$ makes the moment diverge;
the last such power gives logarithmic divergence. For complex coefficients,
one real component has the same nonzero leading term. Conversely,
\eqref{stdil:zeros} makes the weighted tail $O(x^{-2})$, while the
zero endpoint is $O(x^{m-1}(1+|\log x|^{q+1}))$.
The convergent integral is holomorphic in a strip containing $u=m$.
Its removable value in \eqref{stdil:mellin} is
\eqref{stdil:moment}, since $Z(1-m)=0$ and
$1-u=1-m-(u-m)$. Finally \eqref{stdil:subtraction} follows directly
by multiplying the defining scalar Laurent germs by $c_j/h_j$.
\end{proof}

\subsection{A universal construction and elementary moment values}

Fix $\lambda>1$, $q\ge0$, and $r\ge1$, and define
\begin{align}
 P_{q,r}(y)&=\prod_{\ell=q}^{q+r-1}(1-\lambda^{-\ell}y)
            =\sum_{j=0}^rp_jy^j,\notag\\
 H_{q,r}^{\lambda}(x,a)&=
       \sum_{j=0}^r p_j\lambda^{-j}C_q(x/\lambda^j,a).
       \label{stdil:geometric}
\end{align}
The parameters $p_j$ depend on $q,r,\lambda$ and are independent of $a$.

\begin{theorem}[Arbitrarily long ordinary moment blocks]
\label{stdil:block}
For every $a>0$, the kernel in \eqref{stdil:geometric} is
$O(x^{-q-r-1})$ at infinity, so its Mellin integral converges absolutely
on $0<\Re u<q+r+1$. For every $k=q,\ldots,q+r-1$,
\begin{align}
 \int_0^\infty x^kH_{q,r}^{\lambda}(x,a)\,dx
  &={(-1)^{k+1}\over k+1}Q_q(k+1)B_{k+1}(a)\log\lambda
       \prod_{\substack{\ell=q\\\ell\ne k}}^{q+r-1}
                  (1-\lambda^{k-\ell}).\label{stdil:explicit}
\end{align}
The empty product is one. If $B_{k+1}(a)=0$, the displayed ordinary
moment is zero: the Mellin numerator has a double zero there.
No derivative of the Hurwitz zeta function occurs in these values.
\end{theorem}
\begin{proof}
In Theorem~\ref{stdil:criterion}, take
$D(s)=P_{q,r}(\lambda^{-s})$. Then $D(-k)=0$ throughout the
specified block. Its roots are simple because $\lambda>1$, and
\[
 D'(-k)=\log\lambda
       \prod_{\substack{\ell=q\\\ell\ne k}}^{q+r-1}
                       (1-\lambda^{k-\ell}).
\]
Thus $Z'(-k)=D'(-k)\zeta(-k,a)$, and
\eqref{stdil:explicit} follows from \eqref{stdil:moment}.
The tail estimate follows from \eqref{stdil:tail}. When $q=0$ the
factor with $\ell=0$ is simply $1-y$; no division by $1-\lambda^0$
is made anywhere in the argument.
\end{proof}

\begin{proposition}[Minimal number of scales for uniform cancellation]
\label{stdil:minimal}
Suppose coefficients and distinct positive scales, independent of $a$,
give the decay $H_q(x)=O(x^{-q-r-1})$ for every $a$ in a nonempty
open interval of positive shifts. A nonzero such combination requires
at least $r+1$ scales. For the prescribed geometric scales
$1,\lambda,\ldots,\lambda^r$, the choice
\eqref{stdil:geometric} is unique under the normalization $c_0=1$.
\end{proposition}
\begin{proof}
Each polynomial $B_{k+1}(a)$ is nonzero as a polynomial, so uniform
cancellation forces $\sum_jc_jh_j^k=0$ for $q\le k\le q+r-1$.
For at most $r$ scales the first corresponding square submatrix is
a Vandermonde matrix multiplied by the nonzero factors $h_j^q$.
Consequently all coefficients vanish. On $r+1$ prescribed scales the
nullspace is one-dimensional. For geometric scales its coefficient
polynomial must be a multiple of $P_{q,r}$, and the constant coefficient
normalizes that multiple. The claim concerns uniform cancellation:
special fixed shifts with Bernoulli zeros can need fewer scales.
\end{proof}

For $r=1$ the result is especially transparent:
\begin{equation}\label{stdil:one}
 \boxed{\displaystyle
 \int_0^\infty x^q
  \left[C_q(x,a)-\lambda^{-q-1}C_q(x/\lambda,a)\right]dx
     ={(-1)^{q+1}B_{q+1}(a)\over q+1}\log\lambda.}
\end{equation}
This also has an elementary cutoff proof: if
$f(x)=x^qC_q(x,a)$, then $f(x)=A/x+O(x^{-2})$ with
$A=(-1)^q\zeta(-q,a)$, and the cutoff difference of integrals is
$\int_{R/\lambda}^{R}f(x)\,dx\longrightarrow A\log\lambda$.
The full combined integrand is absolutely integrable. This is the
usual cancellation mechanism behind Frullani-type identities.

For concrete examples put $G(x)=C_0(x,1)=\log x-\psi(x+1)$ and
\[
 U(x)=G(x)-\tfrac34G(x/2)+\tfrac18G(x/4).
\]
At $q=0,r=2,\lambda=2$, the first two moments are
\begin{equation}\label{stdil:digamma-pair}
 \int_0^\infty U(x)\,dx=-\frac{\log2}4,\qquad
 \int_0^\infty xU(x)\,dx=-\frac{\log2}{12}.
\end{equation}
There is a third ordinary moment because $\zeta(-2)=0$:
\begin{equation}\label{stdil:digamma-third}
 \int_0^\infty x^2U(x)\,dx
     =3\zeta'(-2)=-\frac{3\zeta(3)}{4\pi^2}.
\end{equation}
The next tail is nonzero, so the weight $x^3$ already diverges.
As a first generalized-Stieltjes example, put
\[
 V(x)=C_1(x,1)-\tfrac38C_1(x/2,1)+\tfrac1{32}C_1(x/4,1),
 \qquad C_1(x,1)=\tfrac12\log^2x-\psi(x+1)\log x-\gamma_1(x+1).
\]
Then
\begin{equation}\label{stdil:first-index}
 \int_0^\infty xV(x)\,dx=\frac{\log2}{24},\qquad
 \int_0^\infty x^2V(x)\,dx=0.
\end{equation}
The second identity uses both $D(-2)=0$ and $\zeta(-2)=0$.

\subsection{The principal-character completion}

Let $\chi_0$ be principal modulo $d$, put
$E_d(s)=\prod_{p\mid d}(1-p^{-s})$, and let $d_0$ be the product
of the distinct prime divisors of $d$. Empty products at $d=1$ are one.
The correct balanced principal kernel is
\begin{align}
 K_{q,\chi_0}(x)
   &=\frac1d\sum_{\substack{1\le r\le d\\(r,d)=1}}
           C_q(x/d,r/d)\notag\\
   &=\sum_{h\mid d_0}\frac{\mu(h)}h C_q(x/h,1).
      \label{stdil:principal-kernel}
\end{align}
Indeed inclusion--exclusion in the Hurwitz series proves equality
of the Laurent germs. Both have scalar subtraction
$E_d(1)/z=\varphi(d)/(dz)$, which is essential here.

\begin{proposition}[All principal-character extra integer moments]
\label{stdil:principal}
For every $d\ge1$ and every even $q=2n\ge2$, precisely one extra
ordinary integer moment exists, and
\begin{equation}\label{stdil:principal-even}
 \int_0^\infty x^{2n}K_{2n,\chi_0}(x)\,dx
  =E_d(-2n)\frac{(-1)^n(2n)!}{2(2\pi)^{2n}}\zeta(2n+1).
\end{equation}
There are no extra ordinary integer moments at odd orders.
At order zero there is no extra moment for $d=1$. For $d>1$
there is precisely one, with value
\begin{equation}\label{stdil:principal-zero}
 \int_0^\infty K_{0,\chi_0}(x)\,dx
   =-\tfrac12 E_d'(0)
   =\begin{cases}-\tfrac12\log p,&d\text{ is a positive power of a prime }p,\\
                  0,&d\text{ has at least two distinct prime divisors}.
     \end{cases}
\end{equation}
All asserted convergent integrals are ordinary and absolutely convergent.
\end{proposition}
\begin{proof}
Apply Theorem~\ref{stdil:criterion} with $D=E_d$ and $a=1$.
For $k\ge1$ none of the Euler factors $1-p^k$ vanishes. The simple
zeros of $\zeta$ at the negative even integers, its nonzero values
at negative odd integers, and
\[
 \zeta'(-2n)=\frac{(-1)^n(2n)!}{2(2\pi)^{2n}}\zeta(2n+1)
\]
are the classical consequences of its functional equation
\cite{stchar:classical}. They prove all assertions for $q\ge1$,
including failure of the next moment. At zero $\zeta(0)=-1/2$;
$E_d$ has a zero whose multiplicity is the number of distinct primes
dividing $d$, while $E_d(-1)\zeta(-1)\ne0$. These facts give
\eqref{stdil:principal-zero}, again with no longer moment block.
\end{proof}
